\section{Conclusão}
\label{sec:conclusao}

Esta revisão sistemática mapeou 28 estudos sobre persistência de
conhecimento em agentes de codificação baseados em IA, todos
publicados a partir de 2024.
O campo é recente e predominantemente composto por preprints (82\%),
mas já apresenta convergência em torno de padrões arquiteturais e
de avaliação recorrentes.
A persistência de memória produz ganhos em todos os cenários
testados, porém com rendimentos decrescentes, riscos de degradação
documentados e lacunas persistentes na documentação de custos.

As cinco questões secundárias receberam respostas com diferentes
graus de certeza.
Quanto à taxonomia (QS1), bancos de experiência predominam (46\%),
o LLM atua como controlador em 54\% dos sistemas, e o escopo
temporal híbrido é o mais frequente (58\%).
Sobre avaliação (QS2), SWE-bench Verified concentra 67\% dos
estudos com benchmark, porém a comparabilidade é limitada por
diferenças de modelo, framework e condições experimentais.
Os ganhos de desempenho (QS3) são universalmente positivos nas 19
comparações controladas, com mediana de 6,8~pp e amplitude de
+2,2 a +39,0~pp, e rendimentos decrescentes em baselines acima
de 65\%.
Em relação ao custo (QS4), 43\% dos estudos não reportam dados de
custo, e o impacto da memória é bimodal: sistemas que substituem
contexto reduzem tokens, enquanto sistemas que adicionam contexto
elevam o custo em aproximadamente 5 a 21\%.
Sobre complexidade e eficácia (QS5), 32\% dos estudos reportam
degradação por memória, com ruído na recuperação como mecanismo
principal, e não há evidência de que arquiteturas mais complexas
produzam ganhos sistematicamente superiores.

Nenhum dos 28 estudos compara arquiteturas de persistência
diretamente.
Essa ausência, somada à subnotificação de custos e à incerteza
sobre o papel da complexidade arquitetural, impede decisões
informadas sobre qual abordagem adotar em cada contexto.
A segunda etapa desta dissertação endereçará a primeira dessas
lacunas por meio de um experimento comparativo com ao menos três
arquiteturas representativas, sob o mesmo benchmark, o mesmo
modelo e com análise integrada de custo por tarefa.
